\subsection{Censo completo de cilindros y selección de la frontera terminal}
\label{censo:terminal-completo}

La lectura terminal utiliza conjuntamente dos informaciones: la compatibilidad
de los cilindros regionales y un perfil de incidencia con canales y posiciones
etiquetados. En este apartado se construyen los catálogos completos, se
enumeran sus ternas y se demuestra el efecto de cada condición del lector.
El punto de partida son las salidas regionales ya producidas en la
sección~\ref{sec:generacion}, no valores introducidos para definir el
transporte APP--TRIT--TPK. La reducción a una matriz pertenece a una lectura
posterior de esas salidas y conserva el orden de los tres canales.

El resultado es un teorema censal con estructura de partida explícita.
El perfil de Gram, pesos, distancias y márgenes se declara como dato del
lector terminal. La exhaustividad de la enumeración demuestra qué selecciona
ese perfil sobre los catálogos fijados; no demuestra, por sí sola, que el
perfil sea una ley universal ni que pueda reconstruirse desde la única terna
que finalmente lo satisface. Esta separación permite incorporar la prueba
finita completa sin invertir su genealogía.

\subsubsection{De las salidas regionales a los catálogos ternarios}

Sean \(x_\pi,x_e,x_\varphi\in[0,1)\) las partes fraccionarias de las
tres coordenadas regionales ya construidas. Se utilizan sus cilindros
publicados de mil posiciones decimales,
\[
 J_\chi^{(1000)}=
 \left[\frac{d_\chi}{10^{1000}},
       \frac{d_\chi+1}{10^{1000}}\right),
 \qquad \chi\in\{\pi,e,\varphi\}.
\]
Los enteros \(d_\chi\) son lecturas finitas de salida. La longitud mil
es suficiente para los enteros que siguen, como se comprueba mediante las
dos cotas de cada intervalo; no es una profundidad supuesta infinita.
Definimos
\begin{equation}
 p=30,\quad r=1050,\quad L=p+r=1080,\quad k=171,
 \qquad D=3^L,\quad Q=1000^k,
 \label{censo:parametros}
\end{equation}
y
\[
 P_\chi=\lfloor3^{30}x_\chi\rfloor,\qquad
 T_\chi=\lfloor Qx_\chi\rfloor.
\]
Aquí \(k\) cuenta tríadas decimales y no designa el registro dodecafásico
\(K\). Se tiene \(1000^{171}\leq3^{1080}<1000^{172}\).

\begin{lemma}[Extracción estable de un entero]
\label{censo:estabilidad}
Para \(d,m,h\in\mathbb Z\), con \(m,h>0\), el valor de
\(\lfloor mx\rfloor\) es constante en \([d/h,(d+1)/h)\) si y sólo si
\begin{equation}
 \left\lfloor\frac{md}{h}\right\rfloor
 =\left\lfloor\frac{m(d+1)-1}{h}\right\rfloor.
 \label{censo:prueba-estabilidad}
\end{equation}
Cuando se cumple, cualquiera de los dos miembros da ese valor.
\end{lemma}
\begin{proof}
El mínimo de \(mx\) se alcanza en el extremo izquierdo. Su supremo
es \(m(d+1)/h\) y no se alcanza. El mayor entero que puede tomar
\(\lfloor mx\rfloor\) es
\(\lceil m(d+1)/h\rceil-1\), igual al segundo miembro porque
el numerador y el denominador son enteros. La igualdad de los extremos
del intervalo de valores prueba la afirmación.
\end{proof}

La prueba de estabilidad se cumple en los tres canales para
\(m=3^{30},Q,3^{1080}\). La última elección se utilizará únicamente
para una comparación posterior, no para seleccionar el perfil. Las
primeras treinta posiciones recuperadas son precisamente las cinco
etapas de seis trits del transporte regional:
\begin{center}\small
\begin{tabular}{@{}c l r@{}}\toprule
Canal&Prefijo de treinta trits&\(P_\chi\)\\\midrule
\(\pi\)&\texttt{010211012222010211002111110221}&29152671743887\\
\(e\)&\texttt{201101121221102011012222102011}&147887858824447\\
\(\varphi\)&\texttt{121200112202121020010210010200}&127247717616687\\\bottomrule
\end{tabular}
\end{center}
La coincidencia conserva la prolongación ya construida; la lectura decimal
no se usa para sustituir retroactivamente sus matrices de transporte.

Fijado un canal, una cola de \(r\) trits se representa por
\(0\leq R<3^r\). Su cilindro y el cilindro de las primeras \(k\)
tríadas son
\begin{equation}
 I_{\chi,R}=\left[\frac{P_\chi3^r+R}{D},
                       \frac{P_\chi3^r+R+1}{D}\right),\qquad
 J_{\chi,T}=\left[\frac{T_\chi}{Q},\frac{T_\chi+1}{Q}\right).
 \label{censo:dos-cilindros}
\end{equation}
El catálogo incluye todas las colas para las que estos dos cilindros
tienen intersección no vacía.

\begin{proposition}[Catálogo por intersección exacta]
\label{censo:catalogos}
Las colas compatibles son los enteros del intervalo \([a_\chi,b_\chi]\),
con
\begin{align}
 a_\chi&=\max\left(0,
       \left\lfloor\frac{T_\chi D}{Q}\right\rfloor-P_\chi3^r\right),
       \label{censo:extremo-a}\\
 b_\chi&=\min\left(3^r-1,
       \left\lfloor\frac{(T_\chi+1)D-1}{Q}\right\rfloor-P_\chi3^r\right).
       \label{censo:extremo-b}
\end{align}
Para los cilindros regionales fijados, los recortes no actúan y se obtiene
\begin{equation}
\begin{array}{c|r|r|r}
 \chi&a_\chi\bmod729&b_\chi\bmod729&b_\chi-a_\chi+1\\\hline
 \pi&67&262&196\\
 e&584&51&197\\
 \varphi&117&312&196
\end{array}
\label{censo:tabla-catalogos}
\end{equation}
\end{proposition}
\begin{proof}
Escribamos \(N=P_\chi3^r+R\). Dos intervalos semiabiertos de
\eqref{censo:dos-cilindros} se cortan exactamente cuando
\[
 NQ<(T_\chi+1)D,\qquad (N+1)Q>T_\chi D.
\]
La segunda desigualdad equivale a
\(N\geq\lfloor T_\chi D/Q\rfloor\); la primera, a
\(N\leq\lfloor((T_\chi+1)D-1)/Q\rfloor\). Restar
\(P_\chi3^r\) y conservar \(0\leq R<3^r\) prueba las fórmulas.
La evaluación por divisiones enteras de los \(T_\chi\) extraídos por
\eqref{censo:prueba-estabilidad} da la tabla. Como \(r\geq6\),
\(P_\chi3^r\equiv0\pmod{729}\), de modo que el residuo del primer
extremo también se obtiene directamente de
\(\lfloor T_\chi D/Q\rfloor\bmod729\).
\end{proof}

Los últimos seis trits de una cola forman la palabra
\(\nu_6(R\bmod729)\), donde \(\nu_6\) es la escritura ternaria
de longitud seis, incluidos sus ceros iniciales. Por la tabla, cada
catálogo tiene menos de \(729\) elementos consecutivos, de manera que
esta reducción es inyectiva sobre él. Así, los tres catálogos de frontera
quedan descritos completamente, sin elegir las ternas finales:
\begin{equation}
 \begin{split}
 \mathcal B_\pi&=\{\nu_6(67+i):0\leq i<196\},\\
 \mathcal B_e&=\{\nu_6((584+j)\bmod729):0\leq j<197\},\\
 \mathcal B_\varphi&=\{\nu_6(117+\ell):0\leq\ell<196\}.
 \end{split}
 \label{censo:fronteras-explicitas}
\end{equation}
Exigir que el extremo izquierdo de \(I_{\chi,R}\) pertenezca a
\(J_{\chi,T}\) sería otra condición: perdería el primer cilindro
que corta la frontera y daría \(195,196,195\). El censo utiliza la
intersección de los cilindros completos, de acuerdo con su definición.

\subsubsection{Perfil declarado e indicadores de compatibilidad}

Para \(u,v\in\mathbb F_3^6\), sean
\(g(u,v)=\sum u_iv_i\in\mathbb F_3\),
\(n(u)=g(u,u)\), \(w(u)=\#\{i:u_i\ne0\}\) y
\(h(u,v)=\#\{i:u_i\ne v_i\}\). Las dos últimas funciones toman
valores enteros; las dos primeras, valores en \(\mathbb F_3\).
El perfil local del lector, en orden \((\pi,e,\varphi)\), es
\begin{equation}
 \begin{split}
 (g_{\pi e},g_{\pi\varphi},g_{e\varphi})&=(2,2,0),\qquad
 (n_\pi,n_e,n_\varphi)=(0,0,0),\\
 (w_\pi,w_e,w_\varphi)&=(3,3,3),\qquad
 (h_{\pi e},h_{\pi\varphi},h_{e\varphi})=(2,5,3).
 \end{split}
 \label{censo:perfil-local}
\end{equation}
Las palabras se ordenan como en \eqref{censo:fronteras-explicitas}.
Usaremos \([E]\in\{0,1\}\) para el indicador de una condición \(E\).
Para cada par de canales \(\chi,\psi\), la matriz rectangular
\(A_{\chi\psi}\) tiene entrada \([g(u,v)=g_{\chi\psi}]\).
La matriz \(A^{h}_{\chi\psi}\) multiplica esa entrada por
\([h(u,v)=h_{\chi\psi}]\). Sean \(D_\chi^n\) la matriz diagonal
de \([n(u)=0]\), y \(D_\chi^{nw}\) la de
\([n(u)=0][w(u)=3]\).

\begin{theorem}[Enumeración exhaustiva por sumas finitas]
\label{censo:teorema-conteo}
El producto de los tres catálogos contiene \(7\,567\,952\) ternas.
Las condiciones sucesivas de \eqref{censo:perfil-local} dejan
\begin{equation}
 7\,567\,952\longrightarrow275\,794\longrightarrow6235
 \longrightarrow3127\longrightarrow331.
 \label{censo:cadena-local}
\end{equation}
Estos cuatro conteos son, respectivamente,
\begin{align}
 N_g&=\operatorname{tr}(A_{\pi e}A_{e\varphi}A_{\varphi\pi}),
 \nonumber\\
 N_{gn}&=\operatorname{tr}(D_\pi^n A_{\pi e}
             D_e^n A_{e\varphi}D_\varphi^n A_{\varphi\pi}),
 \nonumber\\
 N_{gnw}&=\operatorname{tr}(D_\pi^{nw} A_{\pi e}
             D_e^{nw} A_{e\varphi}D_\varphi^{nw} A_{\varphi\pi}),
 \nonumber\\
 N_{gnwh}&=\operatorname{tr}(D_\pi^{nw} A^h_{\pi e}
             D_e^{nw} A^h_{e\varphi}D_\varphi^{nw} A^h_{\varphi\pi}).
 \label{censo:trazas}
\end{align}
Todos los productos y trazas de esta fórmula se calculan en los enteros,
no en el cuerpo de tres elementos.
\end{theorem}
\begin{proof}
El tamaño del producto es \(196\cdot197\cdot196\). Al expandir
una traza de \eqref{censo:trazas}, sus índices recorren exactamente una
terna de palabras. El sumando vale uno si satisface todas las condiciones
representadas y cero en otro caso. Cada terna se cuenta una vez.
Las matrices diagonales añaden las condiciones individuales; los factores
rectangulares añaden las condiciones por pares. Esto demuestra que las
trazas son precisamente los cardinales anunciados.

Para explicitar su evaluación, fijados los índices \(i,j\) de los dos
primeros canales, se forman los conjuntos de índices \(\ell\) del tercero
que cumplen cada condición por pares. Se intersectan sus indicadores,
se añaden los indicadores individuales pertinentes y se suma su cardinal.
La suma final recorre \(0\leq i<196\), \(0\leq j<197\).
La sustitución de las palabras de \eqref{censo:fronteras-explicitas} y
del perfil \eqref{censo:perfil-local} da, en ese orden,
\(275794,6235,3127,331\). El suplemento registra las cuatro contribuciones
de cada uno de los \(196\) índices \(i\) y las \(331\) ternas completas;
su programa evalúa estas sumas con indicadores binarios exactos y contrasta
escalarmente cada superviviente. No interviene una muestra ni una tolerancia.
\end{proof}

\subsubsection{Incidencia etiquetada y carga de orientación}

Para una terna superviviente, sea \(B\in\mathbb F_3^{3\times6}\) su
matriz de filas ordenadas. El segundo perfil exige
\begin{equation}
 \operatorname{rank}_{\mathbb F_3}B=3,\qquad
 w_{\mathrm{col}}(B)=(1,1,2,2,3,0),\qquad
 q_{\partial}(B)=(1,2,2,1,2,0),
 \label{censo:perfil-columnas}
\end{equation}
donde \(w_{\mathrm{col}}\) cuenta los elementos no nulos de cada
columna y \(q_{\partial}\) suma sus entradas módulo tres. Es un perfil
de columnas etiquetadas; permutarlas sin transportar el perfil define otro
lector. Denotemos por \(\mathcal F\) el conjunto de las \(331\) ternas
del teorema anterior.

\begin{proposition}[Reducción del censo al par terminal]
\label{censo:dos-matrices}
Las sumas de indicadores sobre \(\mathcal F\), al exigir sucesivamente
rango, pesos de columna y sumas de columna, son
\begin{equation}
 \sum_{B\in\mathcal F}[\operatorname{rank}B=3]=331,
 \qquad N_{\mathrm{rango,pesos}}=18,
 \qquad N_{\mathrm{rango,pesos,sumas}}=2.
 \label{censo:conteo-columnas}
\end{equation}
Con índices de catálogo comenzando en uno, las dos ternas finales son
\((120,197,157)\) y \((129,197,148)\), y sus matrices son
\[
 B_0=\begin{pmatrix}0&2&0&2&2&0\\0&0&1&2&2&0\\1&0&1&0&1&0\end{pmatrix},
 \qquad
 B_1=\begin{pmatrix}0&2&1&0&2&0\\0&0&1&2&2&0\\1&0&0&2&1&0\end{pmatrix}.
\]
\end{proposition}
\begin{proof}
En cada terna del conjunto ya enumerado se calcula el rango por eliminación
en \(\mathbb F_3\), se cuentan los elementos no nulos de las seis
columnas y se suman éstas módulo tres. Los productos de los indicadores
de \eqref{censo:perfil-columnas} dan \eqref{censo:conteo-columnas}.
La sustitución de los dos índices finales en
\eqref{censo:fronteras-explicitas} da las matrices expuestas. La lista
completa de supervivientes y los filtros exactos del suplemento permiten
reproducir la eliminación sin presuponer esas dos matrices como dominio.
\end{proof}

La orientación \(\sigma=(1,1,-1)=(1,1,2)\in\mathbb F_3^3\)
selecciona la recta
\(\langle\sigma\rangle=\{(0,0,0),(1,1,2),(2,2,1)\}\).
Las cargas de filas de \(B_0,B_1\) son \((0,2,0)\) y \((2,2,1)\);
sólo la segunda pertenece a esa recta. Se recupera así el último paso
del selector de la sección~\ref{subsec:k-terminal}, ahora precedido por
la construcción y enumeración de su dominio completo. La lectura estable
de \(\lfloor3^{1080}x_\chi\rfloor\), efectuada después de los filtros,
da las fronteras \texttt{021020}, \texttt{001220}, \texttt{100210},
con los mismos rangos \((129,197,148)\). Este cotejo posterior no define
el perfil ni participa en las sumas de indicadores.

\subsubsection{Información conservada y alcance del resultado}

El procedimiento conserva el prefijo de treinta trits, la cola compatible,
el orden regional, la posición de cada trit y la orientación. La reducción
módulo \(729\) conserva aquí la identidad de cada candidato porque el
intervalo de colas tiene longitud menor que \(729\); esa inyectividad
no se afirma para una historia arbitraria. El censo de \(331\) demuestra
también que el perfil local aislado no selecciona una sola matriz: los
márgenes etiquetados y la carga aportan condiciones efectivas adicionales.

Se ha recuperado, por tanto, una enumeración finita íntegra y su selector
sobre cilindros de salida, con perfil de lectura explícito. La construcción
del registro firmado utiliza además su lector de memoria y sus canales,
como se especifica a continuación; la matriz visible no se identifica con
ese registro por igualdad de número de componentes. Tampoco este conteo
demuestra la universalidad del perfil, genera de nuevo las coordenadas
regionales ni decide una afirmación sobre alfa o sobre el continuo completo.
